\section{Mô hình hóa bài toán và ngôn ngữ manh mối}

\subsection{Biến chính và tri thức công khai}

Trên bàn $N\times N$, với $N\in\{3,4,5\}$, cột A có chỉ số $c=1$ và hàng trên cùng có $r=1$. Biến $C_i$ mang giá trị True (Criminal, tội phạm) hoặc False (Innocent, vô tội), với chỉ số $i=(r-1)N+c$ chạy từ 1 đến $N^2$ (chẳng hạn A3 trên bàn $3\times3$ là biến số 7). Tên và nghề nghiệp chỉ dùng để hiển thị nhận diện nhân vật, còn mọi lời khai đã mở đều luôn đúng kể cả khi chủ thẻ là tội phạm.

Gọi $L_t$ là tập lời khai đã mở và $S_t$ là tập ô đã biết ở bước $t$. Với mỗi ô trong $S_t$, đặt $\ell_i=C_i$ nếu ô là tội phạm và $\ell_i=\neg C_i$ nếu ô vô tội. Cơ sở tri thức (KB) có dạng

\[
KB_t=\bigwedge_{q\in L_t}\operatorname{CNF}(q)\ \land\ \bigwedge_{i\in S_t}(\ell_i).
\]

Tác tử duyệt các ô chưa mở theo thứ tự hàng, chọn ô đầu tiên có phán quyết bắt buộc. Khi engine chấp nhận, trạng thái và lời khai của ô cùng gia nhập KB ở bước kế tiếp. Vòng lặp dừng khi mở hết bàn, hết phán quyết chứng minh được hoặc gặp mâu thuẫn.

\begin{table}[H]
\centering\small\singlespacing
\caption{Sáu kiểu manh mối cơ sở}
\smallskip
\renewcommand{\arraystretch}{1.23}
\begin{tabularx}{\linewidth}{@{}l >{\raggedright\arraybackslash}X@{}}
\toprule
\textbf{Kiểu} & \textbf{Nội dung trên một phép gán} \\
\midrule
\texttt{FACT} & Một người có trạng thái chỉ định, biểu diễn bằng $C_i$ hoặc $\neg C_i$. \\
\texttt{SAME} & Hai người cùng trạng thái, $C_i\leftrightarrow C_j$. \\
\texttt{DIFFERENT} & Hai người khác trạng thái, $C_i\oplus C_j$. \\
\texttt{EXACTLY} & Số tội phạm trong vùng $R$ đúng bằng $k$. \\
\texttt{AT\_LEAST} & Có ít nhất $k$ tội phạm trong $R$. \\
\texttt{AT\_MOST} & Có nhiều nhất $k$ tội phạm trong $R$. \\
\bottomrule
\end{tabularx}
\end{table}

Các manh mối đếm yêu cầu $k$ nguyên và $0\le k\le |R|$. Câu “không có tội phạm” tương ứng với $\operatorname{EXACTLY}(0,R)$. Với vùng rỗng, số tội phạm bằng 0.

\clearpage

\subsection{Vùng, so sánh và điều kiện lồng}

\begin{table}[H]
\centering\small\singlespacing
\caption{Vùng cơ sở và các phép hợp thành}
\smallskip
\renewcommand{\arraystretch}{1.23}
\begin{tabularx}{\linewidth}{@{}l >{\raggedright\arraybackslash}X@{}}
\toprule
\textbf{Kiểu vùng} & \textbf{Tập ô được chọn} \\
\midrule
\texttt{ROW}, \texttt{COLUMN} & Toàn bộ một hàng hoặc một cột hợp lệ. \\
\texttt{NEIGHBORS} & Các ô kề theo tám hướng, bỏ ô mốc. Một ô góc có ba láng giềng, ô biên còn lại có năm và ô ở trong có tám. \\
\texttt{EXPLICIT} & Các ô phân biệt được liệt kê theo định danh. Danh sách có thể rỗng. \\
\texttt{CORNERS} & Bốn ô góc của bàn. \\
\texttt{UNION}, \texttt{INTERSECTION} & Hợp hoặc giao của hai vùng con. Mỗi ô chỉ xuất hiện một lần trong vùng kết quả. \\
\texttt{DIFFERENCE} & Các ô thuộc vùng trái và không thuộc vùng phải. Thứ tự hai toán hạng quyết định kết quả. \\
\bottomrule
\end{tabularx}
\end{table}

\texttt{COMPARISON} so sánh số tội phạm trong hai vùng bằng một trong sáu quan hệ $=,\ne,<,\le,>,\ge$. Nếu hai vùng giao nhau, cùng một biến ô tham gia vào cả hai tổng. \texttt{CONDITIONAL} nhận hai manh mối bất kỳ làm tiền đề và kết luận. Hai vế có thể tiếp tục chứa điều kiện lồng. Mệnh đề kéo theo sai đúng trong trường hợp tiền đề đúng và kết luận sai.

Trong màn \texttt{case-3x3-01}, lời khai công khai B1 đòi hỏi ít nhất ba tội phạm trong tập láng giềng của A2 sau khi trừ hàng 2. Trên bàn này, vùng còn lại là $\{\mathrm{A1},\mathrm{B1},\mathrm{A3},\mathrm{B3}\}$. Cách phân giải vùng xác định chính xác bốn biến đi vào bộ đếm và bốn ô giao diện cần tô sáng.

Bộ đăng ký ánh xạ tên kiểu JSON sang lớp xử lý. \texttt{@register\_region} gắn kiểu vùng với hàm phân giải, còn \texttt{@register\_clue} gắn kiểu manh mối với kiểm tra tham số, biểu thức logic và hàm đánh giá. Gốc có độ sâu 0. Bộ kiểm tra từ chối cây vùng hoặc cây điều kiện khi độ sâu vượt 12. Nhờ giao diện chung, thêm luật mới chỉ cần đăng ký lớp xử lý và kiểm thử ngữ nghĩa của luật đó.
